\begin{idstatement}{D.2.4}{Proved sign in the full normal-crossings class}
\label{st:D.2.4}
For the threefolds and maps in D.2 whose reduced fibre supports
all have normal crossings, the Euler number is
\[
c_3(Y)=0,\qquad b_3(Y)=2.
\]
All multiplicities, including unequal multiplicities on a
reducible support, are allowed. In this class there are no
positive-Euler-number examples with trivial monodromy.
\end{idstatement}
\begin{proof}
Statement C.2.NC.5 gives $c_3(Y)=2r$, where $r$ counts the
irreducible fibre supports with rational normalization. A fibre
of that kind is necessarily $mS$ for one positive integer $m$.
Choose a small disc around its value and perform the normalized
base change $t=u^m$. Statement C.2.NC.3 identifies the resulting
map to the original neighbourhood as finite \emph{\'etale} of
degree $m$, and its central fibre $S'$ as reduced with normal
crossings. The total space is therefore smooth. Its nearby
fibres are the original connected tori. Stein factorization
over the normal disc shows that its central fibre is connected
as well: the finite Stein factor has generic degree one and is
the disc itself.

Normalization commutes with this \emph{\'etale} base change.
Each normalized component of $S'$ is an unramified finite cover
of the smooth rational normalization $T$ of $S$. Since $T$ is
simply connected, these normalized components are copies of
$T$, and all have geometric genus zero. The local monodromy
after base change is the $m$th power of the original monodromy,
so remains the identity on $H^1$. This contradicts D.2.3.
Therefore $r=0$, proving $c_3(Y)=0$. Poincar\'e duality and
$b_1(Y)=b_2(Y)=0$ give $c_3(Y)=2-b_3(Y)$, hence $b_3(Y)=2$.
The example of D.2.7 realizes this normal-crossings class.
\end{proof}


% D.2.5 from part_D_sign.tex

